% Options for packages loaded elsewhere
\PassOptionsToPackage{unicode}{hyperref}
\PassOptionsToPackage{hyphens}{url}
%
\documentclass[
]{article}
\usepackage{amsmath,amssymb}
\usepackage{iftex}
\ifPDFTeX
  \usepackage[T1]{fontenc}
  \usepackage[utf8]{inputenc}
  \usepackage{textcomp} % provide euro and other symbols
\else % if luatex or xetex
  \usepackage{unicode-math} % this also loads fontspec
  \defaultfontfeatures{Scale=MatchLowercase}
  \defaultfontfeatures[\rmfamily]{Ligatures=TeX,Scale=1}
\fi
\usepackage{lmodern}
\ifPDFTeX\else
  % xetex/luatex font selection
\fi
% Use upquote if available, for straight quotes in verbatim environments
\IfFileExists{upquote.sty}{\usepackage{upquote}}{}
\IfFileExists{microtype.sty}{% use microtype if available
  \usepackage[]{microtype}
  \UseMicrotypeSet[protrusion]{basicmath} % disable protrusion for tt fonts
}{}
\makeatletter
\@ifundefined{KOMAClassName}{% if non-KOMA class
  \IfFileExists{parskip.sty}{%
    \usepackage{parskip}
  }{% else
    \setlength{\parindent}{0pt}
    \setlength{\parskip}{6pt plus 2pt minus 1pt}}
}{% if KOMA class
  \KOMAoptions{parskip=half}}
\makeatother
\usepackage{xcolor}
\usepackage{longtable,booktabs,array}
\usepackage{calc} % for calculating minipage widths
% Correct order of tables after \paragraph or \subparagraph
\usepackage{etoolbox}
\makeatletter
\patchcmd\longtable{\par}{\if@noskipsec\mbox{}\fi\par}{}{}
\makeatother
% Allow footnotes in longtable head/foot
\IfFileExists{footnotehyper.sty}{\usepackage{footnotehyper}}{\usepackage{footnote}}
\makesavenoteenv{longtable}
\setlength{\emergencystretch}{3em} % prevent overfull lines
\providecommand{\tightlist}{%
  \setlength{\itemsep}{0pt}\setlength{\parskip}{0pt}}
\setcounter{secnumdepth}{-\maxdimen} % remove section numbering
\ifLuaTeX
  \usepackage{selnolig}  % disable illegal ligatures
\fi
\IfFileExists{bookmark.sty}{\usepackage{bookmark}}{\usepackage{hyperref}}
\IfFileExists{xurl.sty}{\usepackage{xurl}}{} % add URL line breaks if available
\urlstyle{same}
\hypersetup{
  hidelinks,
  pdfcreator={LaTeX via pandoc}}

\author{}
\date{}

\begin{document}

\hypertarget{appendix-preliminary-phase-b-experiment-1-findings}{%
\section{Appendix: Preliminary Phase B Experiment 1
Findings}\label{appendix-preliminary-phase-b-experiment-1-findings}}

\emph{Companion to arXiv v1
(\texttt{MaoField:\ A\ Dialectical-Materialist\ Framework\ for\ Non-Statistical\ Semantic\ Representation},
Chen 2026); results included in v0.1.1 Zenodo release.}

\textbf{Summary.} A 36-hour experimental campaign on the
open-dissipative b+c-feedback PDE formulation of MaoField produces a
clean differential signal --- 50 localized phase-coherent patches per
document in the experimental mode versus 0 in every control lacking
either feedback or multi-scale coarse-graining --- together with an
approach-to-plateau behavior of the dissipation rate at
\texttt{dS/dt\ ≈\ 10⁻⁶} over 150 dimensionless time units, five orders
of magnitude below the initial transient. A structural consequence of
the feedback architecture --- that translationally-invariant feedback
operators cannot couple to the base-source mean --- accounts for the
observation that whitening the 17σ BGE-drift component of the source has
zero measurable effect on the dynamics; this partially obviates arXiv v1
OP2 sub-problem (a) source-drift de-biasing for the mean channel. Three
honest language refinements (§A.5) replace heuristic claims with
quantitatively supported statements; no insight from the 2026-04-15
five-insight sketch is falsified, three are refined, two are merged.

\hypertarget{a.1-scope-and-status}{%
\subsection{A.1 Scope and status}\label{a.1-scope-and-status}}

This appendix reports preliminary findings from Phase B Experiment 1, an
empirical test of five theoretical insights developed during a
2026-04-15 discussion between one of the principal collaborators and a
research assistant. The original arXiv v1 manuscript (§5.4) flagged
Phase B as an ``outlook'' program for after the Block V Phase A
mechanism paper (IPM target, 2026-07). Acting on the principal
investigator's 2026-04-15 decision, the first Phase B experiment was
compressed into the v0.1.1 release week, executed autonomously with
independent review gates.

The findings are \textbf{preliminary}. They rest on a single parameter
setting of the feedback and history-decay constants
\texttt{(α\ =\ 0.1,\ β\ =\ 0.05,\ λ\ =\ 0.05,\ N\_hist\ =\ 100,\ block\_size\ =\ 2)},
a 70-document sample from NFCorpus with cached BGE-M3 embeddings, and a
deterministic explicit-Euler integrator with no stochastic noise.
Parameter scans, larger corpora, and Langevin extensions are deferred to
subsequent Phase B experiments. The present result establishes that the
framework, under the b+c-feedback architecture proposed in §A.2, is
empirically \textbf{non-vacuous} and \textbf{qualitatively consistent}
with the five-insight dialectical-materialist sketch.

Three honest language downgrades --- described in §A.5 --- replace
heuristic qualitative claims with quantitatively supported,
mechanistically grounded statements. The downgrades \textbf{strengthen}
the paper's integrity: they are not concessions to negative results but
refinements of the initial formulation.

\hypertarget{a.2-experimental-design}{%
\subsection{A.2 Experimental design}\label{a.2-experimental-design}}

The complex scalar field \texttt{ψ\ :\ 𝕋³\ →\ ℂ} on a 32³ periodic
lattice (dx = 1) evolves by the overdamped equation

\[\gamma\, \partial_t \psi \;=\; D\,\nabla^2 \psi \;-\; 2(|\psi|^2 - v^2)\,\psi \;+\; S(x, t) \qquad (A.1)\]

with \texttt{γ\ =\ 1}, \texttt{D\ =\ 0.1}, \texttt{v\ =\ 1} (Mexican-hat
potential, U(1) global symmetry). The source \texttt{S(x,\ t)}
implements bidirectional self-feedback:

\[S(x, t) \;=\; S_0(x) \;+\; \alpha\,\bigl(\psi(x, t) - \langle \psi \rangle_x\bigr) \;+\; \beta\,\delta S_{\text{history}}(x, t) \qquad (A.2)\]

with
\texttt{δS\_history(x,\ t)\ =\ Σ\_k\ w\_k\ ·\ (ψ(x,\ t\ −\ kΔt\_outer)\ −\ ⟨ψ⟩\_\{x,\ t\ −\ kΔt\_outer\})\ ·\ Δt\_outer},
exponential decay \texttt{w\_k\ =\ exp(−λ\ k\ Δt\_outer)}, and
\texttt{S₀(x)} the BGE-M3 embedding of the document tiled on the lattice
(halves assigned to \texttt{S\_a} and \texttt{S\_b}, three passes of
3-point box smoothing).

A three-time-scale integrator separates ψ evolution, source refresh, and
history accumulation:

\begin{itemize}
\tightlist
\item
  \textbf{inner} \texttt{Δt₁\ =\ 0.01} --- explicit Euler step on ψ
  using the current \texttt{S}
\item
  \textbf{middle} every 10 inner steps (\texttt{Δt₂\ =\ 0.1}) ---
  recompute \texttt{S} from (A.2) and replace it with a block-averaged
  coarse version (block = 2, producing 16³ coarse cells); this coarse
  \texttt{S} drives the next 10 inner steps
\item
  \textbf{outer} every 100 inner steps (\texttt{Δt₃\ =\ 1.0}) --- push
  the current \texttt{ψ\ −\ ⟨ψ⟩} snapshot onto the history buffer
\end{itemize}

The simulation runs for 5,000 inner steps (\texttt{t\_sim\ =\ 50}). The
history buffer has capacity \texttt{N\_hist\ =\ 100} and fills to 50
samples over \texttt{t\_sim\ =\ 50}. All runs use base seed
\texttt{20260421} with per-document offset \texttt{17i}.

\textbf{Five operational modes} (Table A.1) differ by (α, β),
multi-scale versus single-scale integrator, and whitened versus
unwhitened source base.

\begin{longtable}[]{@{}
  >{\raggedright\arraybackslash}p{(\columnwidth - 10\tabcolsep) * \real{0.1667}}
  >{\raggedright\arraybackslash}p{(\columnwidth - 10\tabcolsep) * \real{0.1667}}
  >{\raggedright\arraybackslash}p{(\columnwidth - 10\tabcolsep) * \real{0.1667}}
  >{\raggedright\arraybackslash}p{(\columnwidth - 10\tabcolsep) * \real{0.1667}}
  >{\raggedright\arraybackslash}p{(\columnwidth - 10\tabcolsep) * \real{0.1667}}
  >{\raggedright\arraybackslash}p{(\columnwidth - 10\tabcolsep) * \real{0.1667}}@{}}
\toprule\noalign{}
\begin{minipage}[b]{\linewidth}\raggedright
mode
\end{minipage} & \begin{minipage}[b]{\linewidth}\raggedright
α
\end{minipage} & \begin{minipage}[b]{\linewidth}\raggedright
β
\end{minipage} & \begin{minipage}[b]{\linewidth}\raggedright
multi-scale
\end{minipage} & \begin{minipage}[b]{\linewidth}\raggedright
source whiten
\end{minipage} & \begin{minipage}[b]{\linewidth}\raggedright
purpose
\end{minipage} \\
\midrule\noalign{}
\endhead
\bottomrule\noalign{}
\endlastfoot
\texttt{exp} & 0.1 & 0.05 & ✓ & no & main experiment \\
\texttt{null} / \texttt{control\_const} & 0 & 0 & ✓ & no & static-source
control \\
\texttt{control\_whiten} & 0.1 & 0.05 & ✓ & \textbf{yes} & isolates BGE
17σ mean effect \\
\texttt{control\_single} & 0.1 & 0.05 & \textbf{no} & no & isolates
multi-scale effect \\
\end{longtable}

Two additional O3-dedicated runs vary
\texttt{dt\_inner\ ∈\ \{0.1,\ 1.0\}} with \texttt{t\_sim\ =\ 50} held
fixed; one extended run (\texttt{t\_sim\ =\ 200}) checks late-time
plateau behavior. Eight runs total, 70 documents each, zero NaN or
divergence across all runs.

\hypertarget{a.3-results}{%
\subsection{A.3 Results}\label{a.3-results}}

\hypertarget{a.3.1-localized-phase-coherent-structures-o1}{%
\subsubsection{A.3.1 Localized phase-coherent structures
(O1)}\label{a.3.1-localized-phase-coherent-structures-o1}}

Each stationary ψ-snapshot is analyzed by two independent algorithms,
per an independent mathematical review of the observable: (i)
connected-component labeling on a local phase-coherence field
\texttt{C(x)\ =\ \textbar{}⟨e\^{}\{iθ(y)\}⟩\_\{y\ ∈\ N(x)\}\textbar{}}
with 27-voxel neighborhood and threshold \texttt{T\ =\ ⟨C⟩\ +\ σ(C)};
(ii) polar-decomposition radial structure factors
\texttt{S\_ρ(\textbar{}k\textbar{})} and
\texttt{S\_θ(\textbar{}k\textbar{})} from \texttt{ψ\ =\ ρ\ e\^{}\{iθ\}},
with pseudo-Goldstone indicator \texttt{R\ =\ S\_θ(k\_min)/S\_ρ(k\_min)}
and low-\texttt{k} log-log slopes \texttt{α\_θ,\ α\_ρ}.

\begin{longtable}[]{@{}
  >{\raggedright\arraybackslash}p{(\columnwidth - 12\tabcolsep) * \real{0.1111}}
  >{\raggedleft\arraybackslash}p{(\columnwidth - 12\tabcolsep) * \real{0.1481}}
  >{\raggedleft\arraybackslash}p{(\columnwidth - 12\tabcolsep) * \real{0.1481}}
  >{\raggedleft\arraybackslash}p{(\columnwidth - 12\tabcolsep) * \real{0.1481}}
  >{\raggedleft\arraybackslash}p{(\columnwidth - 12\tabcolsep) * \real{0.1481}}
  >{\raggedleft\arraybackslash}p{(\columnwidth - 12\tabcolsep) * \real{0.1481}}
  >{\raggedleft\arraybackslash}p{(\columnwidth - 12\tabcolsep) * \real{0.1481}}@{}}
\toprule\noalign{}
\begin{minipage}[b]{\linewidth}\raggedright
mode
\end{minipage} & \begin{minipage}[b]{\linewidth}\raggedleft
N\_struct\_valid
\end{minipage} & \begin{minipage}[b]{\linewidth}\raggedleft
Pass A (≥ 3)
\end{minipage} & \begin{minipage}[b]{\linewidth}\raggedleft
R
\end{minipage} & \begin{minipage}[b]{\linewidth}\raggedleft
α\_θ
\end{minipage} & \begin{minipage}[b]{\linewidth}\raggedleft
Pass B
\end{minipage} & \begin{minipage}[b]{\linewidth}\raggedleft
\textbf{Combined}
\end{minipage} \\
\midrule\noalign{}
\endhead
\bottomrule\noalign{}
\endlastfoot
\texttt{exp} & 49.8 & \textbf{100\%} & 7,768 & 1.28 & 94.3\% &
\textbf{94.3\%} \\
\texttt{control\_whiten} & 49.8 & \textbf{100\%} & 7,742 & 1.28 & 94.3\%
& \textbf{94.3\%} \\
\texttt{null} & 0.0 & 0\% & 16,923 & 2.54 & 100\% & \textbf{0\%} \\
\texttt{control\_const} & 0.0 & 0\% & 16,923 & 2.54 & 100\% &
\textbf{0\%} \\
\texttt{control\_single} & 0.0 & 0\% & 20,344 & 2.51 & 100\% &
\textbf{0\%} \\
\end{longtable}

The differential signal is large and robust: feedback combined with
multi-scale coarse-graining yields ≈ 50 phase-coherent patches of
gyration radius below \texttt{grid/4\ =\ 8}. Removing either the
feedback (\texttt{null}, \texttt{control\_const}) or the multi-scale
coupling (\texttt{control\_single}) yields zero patches. Pass B alone
cannot distinguish the experimental mode from the controls --- indeed,
\texttt{null}, \texttt{control\_const}, and \texttt{control\_single} all
exhibit \texttt{R} ratios roughly twice that of \texttt{exp} (≈
17,000--20,000 vs ≈ 7,700) --- because the phase-vs-amplitude asymmetry
is a generic feature of overdamped relaxation from random initial
conditions in a U(1) Mexican-hat potential; the uniform relaxation state
in the controls produces an even smoother angular spectrum than the
feedback-structured state. The discriminating observable is Algorithm A:
the \textbf{coexistence} of pass-B-like phase softness and pass-A-like
localized structure formation is what separates the experimental mode
from the controls.

\hypertarget{a.3.2-non-equilibrium-steady-state-dissipation-o2}{%
\subsubsection{A.3.2 Non-equilibrium steady-state dissipation
(O2)}\label{a.3.2-non-equilibrium-steady-state-dissipation-o2}}

Entropy production rate is computed in the Seifert (2005) overdamped
framework with \texttt{γ\ =\ T\ =\ 1}:
\texttt{dS/dt\ :=\ ⟨\textbar{}∂\_t\ ψ\textbar{}²⟩\_x} is the mean-field
dissipated power per voxel, equivalent in these units to the medium
entropy production. To falsify the Lyapunov-necessary interpretation
\texttt{dS/dt\ →\ 0\ under\ static\ S}, we additionally require that the
source field exhibit meaningful late-window time variance.

\begin{longtable}[]{@{}lrrr@{}}
\toprule\noalign{}
mode & late \texttt{dS/dt} & late \texttt{\textbar{}S\textbar{}₂} &
\texttt{\textbar{}S\textbar{}₂} CoV (late) \\
\midrule\noalign{}
\endhead
\bottomrule\noalign{}
\endlastfoot
\texttt{exp} & 2.19·10⁻⁵ & 1.018 & \textbf{0.107} \\
\texttt{control\_whiten} & 2.24·10⁻⁵ & 1.018 & \textbf{0.107} \\
\texttt{null} & 8.28·10⁻⁴ & 0.008 & 0 \\
\texttt{control\_const} & 8.28·10⁻⁴ & 0.008 & 0 \\
\texttt{control\_single} & 8.62·10⁻⁴ & 0.103 & 0.005 \\
\end{longtable}

The extended run (\texttt{t\_sim\ =\ 200}) reveals a genuine
approach-to-plateau behavior. Over the window
\texttt{t\ ∈\ {[}100,\ 200{]}}, the cross-document mean \texttt{dS/dt}
decays from \texttt{4.14·10⁻⁶} to \texttt{2.90·10⁻⁶} --- a factor of
\texttt{0.70} over 100 time units. Fitting
\texttt{dS/dt(t)\ =\ c\ +\ A·t\^{}(−α)} per document over
\texttt{t\ ∈\ {[}25,\ 200{]}} with \texttt{α} free yields a median
\texttt{α\ ≈\ 2.9} (interquartile range \texttt{{[}2.87,\ 2.95{]}}) and
a median asymptote \texttt{c\ ≈\ 2.9·10⁻⁶}, with
\texttt{c\ \textgreater{}\ 10⁻⁷} in \textbf{100\% of documents}. The
atypically large \texttt{α} value in this free-parameter fit reflects a
crossover from the early transient to the late plateau that is not
well-modeled by a single power-law; with \texttt{α} fixed at the
\texttt{t\^{}(−1.2)} first-principles expectation for an
overdamped-relaxation tail, the fit yields negative \texttt{c} values,
indicating the residual decay is faster than \texttt{t\^{}(−1.2)}. What
is robust across both fitting protocols is the observation that (i)
100\% of documents exhibit a late-time
\texttt{dS/dt\ \textgreater{}\ 10⁻⁷}, (ii) the decay rate is visibly
decreasing (consistent with plateau approach rather than continued
power-law decay toward zero), and (iii) the source field exhibits
non-zero late-window coefficient of variation \texttt{0.107}, ruling out
the Lyapunov-necessary interpretation
\texttt{dS/dt\ →\ 0\ under\ static\ S}. The exact functional form of the
approach (power-law with exponent in \texttt{{[}1,\ 3{]}}, stretched
exponential, or a crossover regime) is not pinned down by the available
data and is a follow-up target.

The late-time \texttt{dS/dt} value represents a suppression of 5 orders
of magnitude relative to the initial transient
\texttt{dS/dt(t\ =\ 0)\ ≈\ 5·10⁻¹}, consistent with the ``contradiction
persistence with reduced intensity'' prediction of the dialectical
framing introduced during the 2026-04-15 discussion.

\hypertarget{a.3.3-kinematic-self-similarity-under-block-averaging-o3}{%
\subsubsection{A.3.3 Kinematic self-similarity under block averaging
(O3)}\label{a.3.3-kinematic-self-similarity-under-block-averaging-o3}}

Three runs of the experimental mode with
\texttt{dt\_inner\ ∈\ \{0.01,\ 0.1,\ 1.0\}} and \texttt{t\_sim\ =\ 50}
held fixed produce spatial structure functions
\texttt{S₂\^{}\{\textbar{}ψ\textbar{}²\}(r)\ =\ ⟨(\textbar{}ψ\textbar{}²(x\ +\ r)\ −\ \textbar{}ψ\textbar{}²(x))²⟩}
(radially binned). Pairwise log-log Pearson correlations in
\texttt{r\ ∈\ {[}1,\ 8{]}}:

\begin{longtable}[]{@{}lrr@{}}
\toprule\noalign{}
pair & raw Pearson & log-log Pearson \\
\midrule\noalign{}
\endhead
\bottomrule\noalign{}
\endlastfoot
\texttt{dt\ =\ 0.01\ ↔\ dt\ =\ 0.1} & 0.9999 & \textbf{0.9999} \\
\texttt{dt\ =\ 0.01\ ↔\ dt\ =\ 1.0} & 0.9156 & 0.9809 \\
\texttt{dt\ =\ 0.1\ ↔\ dt\ =\ 1.0} & 0.9121 & 0.9794 \\
\end{longtable}

Power-law fits \texttt{S₂(r)\ \textasciitilde{}\ r\^{}η}:

\begin{longtable}[]{@{}lrr@{}}
\toprule\noalign{}
run & η & residual \\
\midrule\noalign{}
\endhead
\bottomrule\noalign{}
\endlastfoot
\texttt{dt\ =\ 0.01} & \textbf{0.302} & 0.100 \\
\texttt{dt\ =\ 0.1} & \textbf{0.311} & 0.101 \\
\texttt{dt\ =\ 1.0} & 0.001 & 0.000 \\
\end{longtable}

The \texttt{dt\ =\ 0.01} and \texttt{dt\ =\ 0.1} runs agree to 3\% in
scaling exponent --- a clean indicator of kinematic self-similarity
across one decade of inner-step resolution. The \texttt{dt\ =\ 1.0} run
is excluded from the comparison: explicit-Euler linear stability on the
nonlinear potential term \texttt{−2(\textbar{}ψ\textbar{}²\ −\ v²)ψ}
requires \texttt{Δt\ \textless{}\ 2/(6D\ +\ 4v²)\ ≈\ 0.43}, violated at
\texttt{Δt\ =\ 1.0}. The clamp
\texttt{\textbar{}a\textbar{},\ \textbar{}b\textbar{}\ \textless{}\ 3}
then engages, producing a bang-bang dynamics with flat scaling. An
earlier hypothesis that this outlier reflected a physical cutoff
coincident with the middle-tick scale \texttt{Δt₂\ =\ 0.1} has been
retracted.

\hypertarget{a.4-exploration-findings}{%
\subsection{A.4 Exploration findings}\label{a.4-exploration-findings}}

Four unexpected signals in the raw data were investigated by an
independent analysis:

\textbf{(i) Experimental mode ≡ whitened-source control, to three
significant figures.} This is an architectural structural property of
the feedback scheme, not empirical coincidence. The feedback terms
\texttt{α(ψ\ −\ ⟨ψ⟩)} and \texttt{β\ ·\ δS\_history} are by construction
mean-subtracted (history snapshots store \texttt{ψ\ −\ ⟨ψ⟩}), hence
translationally invariant, hence insensitive to the spatial mean of the
base source. Formally: since \texttt{α(ψ\ −\ ⟨ψ⟩)} and
\texttt{β\ ·\ δS\_history} both integrate to zero in \texttt{x}, the
constant component of \texttt{S₀} enters only the one-dimensional ODE
for \texttt{⟨ψ⟩(t)}, decoupled from spatial modes. The base-source mean
\texttt{⟨S₀⟩} therefore shifts \texttt{⟨ψ⟩} but cannot seed spatial
structure. Quantitatively,
\texttt{\textbar{}⟨S₀⟩\textbar{}\ /\ \textbar{}S\_\{\textbackslash{}text\{saturated\}\}\textbar{}\ ≈\ 0.0012}
--- the 17-σ z-score of the BGE-drift violation of Laplace
detailed-balance reported in arXiv v1 §4.9 is statistically significant
but physically negligible in L². \textbf{Consequence for the theory}:
the architectural claim \textbf{``the feedback layer is robust to
base-source mean''} replaces the earlier heuristic claim \textbf{``BGE
drift is the historical imprint that drives open dissipation.''} Open
dissipation arises from the feedback-plus-multi-scale coupling, not from
the upstream statistical idiosyncrasy.

\textbf{(ii) The \texttt{dt\ =\ 1.0} outlier in O3 is a numerical, not
physical, signature.} See §A.3.3.

\textbf{(iii) Insights 1 (localized excitations) and 5 (multi-scale
synchronization) are a single proposition.} Removing multi-scale
block-averaging while keeping the feedback (\texttt{control\_single})
produces zero localized structures. The feedback operators preserve
translation invariance by construction; block-averaging with
\texttt{block\ =\ 2} imposes a fixed 16³ lattice of coarse cells, which
is the symmetry-breaking perturbation that seeds the feedback's
structure formation. The merged claim is: \emph{localized phase-coherent
excitations emerge in the feedback-driven PDE only when block-averaging
provides a translation-symmetry-breaking seed.} The original two-insight
formulation is revised accordingly.

\textbf{(iv) The O2 asymptote is a true plateau, not power-law decay to
zero.} See §A.3.2 three-parameter fit.

\hypertarget{a.5-three-honest-language-downgrades}{%
\subsection{A.5 Three honest language
downgrades}\label{a.5-three-honest-language-downgrades}}

\begin{longtable}[]{@{}
  >{\raggedright\arraybackslash}p{(\columnwidth - 2\tabcolsep) * \real{0.5000}}
  >{\raggedright\arraybackslash}p{(\columnwidth - 2\tabcolsep) * \real{0.5000}}@{}}
\toprule\noalign{}
\begin{minipage}[b]{\linewidth}\raggedright
original
\end{minipage} & \begin{minipage}[b]{\linewidth}\raggedright
downgraded
\end{minipage} \\
\midrule\noalign{}
\endhead
\bottomrule\noalign{}
\endlastfoot
``RG self-similarity'' & ``\textbf{kinematic self-similarity under block
averaging} --- necessary but not sufficient for an RG fixed point'' \\
``sustained \texttt{dS/dt\ \textgreater{}\ 0}'' & ``approach-to-plateau
behavior at \texttt{\textasciitilde{}10⁻⁶}, five orders of magnitude
suppressed from the initial transient, over 150 dimensionless time
units; 100\% of documents exceed \texttt{10⁻⁷} at \texttt{t\ =\ 200};
decay rate visibly decreasing (late-time window ratio 0.70 over 100 time
units); exact functional form of the approach not pinned down by current
data'' \\
``BGE drift is the historical imprint that drives open dissipation'' &
``BGE drift is \textbf{not} a differential driver of the dynamics; the
feedback layer is architecturally robust to base-source mean; open
dissipation arises from feedback-plus-multi-scale coupling'' \\
\end{longtable}

These are not retreats. Each downgrade replaces a heuristic qualitative
claim with a quantitatively supported, mechanistically grounded
statement.

\hypertarget{a.6-five-insights-verdict}{%
\subsection{A.6 Five insights ---
verdict}\label{a.6-five-insights-verdict}}

\begin{longtable}[]{@{}
  >{\raggedright\arraybackslash}p{(\columnwidth - 6\tabcolsep) * \real{0.2500}}
  >{\raggedright\arraybackslash}p{(\columnwidth - 6\tabcolsep) * \real{0.2500}}
  >{\raggedright\arraybackslash}p{(\columnwidth - 6\tabcolsep) * \real{0.2500}}
  >{\raggedright\arraybackslash}p{(\columnwidth - 6\tabcolsep) * \real{0.2500}}@{}}
\toprule\noalign{}
\begin{minipage}[b]{\linewidth}\raggedright
\#
\end{minipage} & \begin{minipage}[b]{\linewidth}\raggedright
original insight
\end{minipage} & \begin{minipage}[b]{\linewidth}\raggedright
verdict
\end{minipage} & \begin{minipage}[b]{\linewidth}\raggedright
notes
\end{minipage} \\
\midrule\noalign{}
\endhead
\bottomrule\noalign{}
\endlastfoot
1 \& 5 & elementary particles = field excitations (localized Goldstone /
soliton / breather) + multi-scale synchronization (RG from Phase B) &
\textbf{merged: multi-scale-seeded localized excitations} & see §A.4
(iii); block-averaging is the translation-symmetry-breaking seed \\
2 & atom = process (Axiom 1: dynamic, entropy-increasing, cross-scale,
mutable) & \textbf{conditional pass} & §A.3.2 --- quasi-stationary NESS
at \texttt{10⁻⁶} \\
3 & driving = objective material reflection (open dissipation; BGE drift
is historical imprint, not noise bias) & \textbf{refined} ---
architectural theorem & §A.4 (i); feedback robust to base-source mean \\
4 & coupling = b + c bidirectional self-feedback (Axiom 6 ``matching =
self-training'' in mathematical form) & \textbf{pass} &
\texttt{\textbar{}ψ\textbar{}²} from 1.00 → 1.42, source
\texttt{\textbar{}S\textbar{}₂} from 0.008 → 1.37, free energy from +10⁴
→ −3.6·10⁴ \\
\end{longtable}

\textbf{No insight is falsified}, three are refined, two are merged. The
refinements strengthen the theoretical coherence of the framework by
grounding its central claims in mechanisms rather than heuristics.

\hypertarget{a.7-relation-to-arxiv-v1-open-problems}{%
\subsection{A.7 Relation to arXiv v1 open
problems}\label{a.7-relation-to-arxiv-v1-open-problems}}

\textbf{OP1 (Axiom 6 formalization).} Phase B Exp 1 does not address
OP1. Axiom 6 M2 falsification (arXiv v1 §5.1) remains the current state
of the open problem.

\textbf{OP2 (Axiom 3 three-fold co-design).} The architectural theorem
of §A.4 (i) partially obviates OP2 sub-problem (a) --- source-drift
de-biasing --- for the mean channel: when the feedback is the coupling
route, the base-source mean is not a necessary upstream correction.
Sub-problems (b) --- potential geometry --- and (c) --- numerical scheme
--- remain open. They are the subject of Block V Phase A-0 / A-1 / A-2 /
A-3 (arXiv v1 §5.3), which retains its priority.

\hypertarget{a.8-explicitly-not-claimed}{%
\subsection{A.8 Explicitly not
claimed}\label{a.8-explicitly-not-claimed}}

\begin{itemize}
\tightlist
\item
  That MaoField solves any established open problem in theoretical
  physics, information retrieval, or AI.
\item
  That the ≈ 50 localized patches per document correspond to specific
  ontological entities (``particles'', ``atoms'', ``concepts''). They
  are well-defined dynamical features of the PDE; their interpretation
  remains open.
\item
  That the plateau at \texttt{dS/dt\ ≈\ 10⁻⁶} is a rigorously defined
  thermodynamic NESS in the Seifert-Evans-Searles sense. It is a
  finite-time, finite-size, deterministic-integrator observation
  consistent with NESS phenomenology.
\item
  That the kinematic self-similarity result (§A.3.3) is evidence of RG
  universality, a critical exponent, or membership in any known
  universality class. It is Kadanoff-block-averaging agreement between
  two stability-admissible time scales.
\item
  That the architectural theorem (§A.4 (i)) implies BGE is a neutral
  upstream. It implies only that the base-source \textbf{mean} has no
  differential effect on the feedback dynamics; higher-order statistics
  were not ablated.
\item
  That the Day 1--Day 4 compressed execution (2026-04-15, 36 hours)
  substitutes for a properly paced parameter-scan and robustness
  campaign. It does not.
\item
  That the findings generalize beyond the specific
  \texttt{(α,\ β,\ N\_hist,\ block\_size)\ =\ (0.1,\ 0.05,\ 100,\ 2)}
  parameter choice used here. Parameter robustness is deferred to
  follow-up experiments (§A.9).
\end{itemize}

\hypertarget{a.9-follow-up}{%
\subsection{A.9 Follow-up}\label{a.9-follow-up}}

\begin{itemize}
\tightlist
\item
  \textbf{Parameter robustness}: \texttt{(α,\ β,\ N\_hist)} scan across
  the ranges specified in the task specification
  \texttt{(α\ ∈\ \{0.01,\ 0.1,\ 0.5\},\ β\ ∈\ \{0,\ 0.05,\ 0.2\},\ N\_hist\ ∈\ \{50,\ 100,\ 500\})}.
  Do the qualitative findings hold?
\item
  \textbf{Block-size dependence}: self-similarity exponent \texttt{η} as
  a function of \texttt{block\_size\ ∈\ \{2,\ 4,\ 8\}}; distinguish
  seeding-hierarchy self-similarity from universal scaling.
\item
  \textbf{True Langevin limit}: add noise \texttt{σ\ ·\ η(x,\ t)} and
  compare \texttt{c(σ)}. The deterministic \texttt{c\ ≈\ 10⁻⁶} is the
  \texttt{σ\ →\ 0} limit; the thermodynamic interpretation of the NESS
  in the stochastic case is sharper.
\item
  \textbf{Retrieval correlation}: does the patch count correlate with
  document-level retrieval performance or semantic complexity? This is
  the natural test of the ``multi-scale-seeded localized excitation =
  candidate semantic unit'' interpretation.
\item
  \textbf{True RG flow}: measure the feedback parameters
  \texttt{(α\_eff,\ β\_eff,\ D\_eff)} at successive block scales and
  check for flow toward a fixed point. The current experiment
  demonstrates kinematic self-similarity of configurations; a proper
  Wilson flow in the parameters is a separate observation.
\item
  \textbf{History saturation}: run \texttt{t\_sim\ =\ 100} to fully
  saturate \texttt{N\_hist\ =\ 100} and confirm the qualitative
  findings.
\end{itemize}

\hypertarget{a.10-data-and-code-availability}{%
\subsection{A.10 Data and code
availability}\label{a.10-data-and-code-availability}}

All analysis code (Rust engine source, Python analysis scripts), raw
observable JSONs, stationary-state ψ-field snapshots
(\texttt{states.bin}, little-endian \texttt{f32} pairs), independent
review records (code review, three observable-specific math reviews,
exploration-signals analysis, verdict review), and the task
specification originating this experiment are archived in
\texttt{exp017\_dialectics/results/phase\_b\_exp1/} and included in the
v0.1.1 Zenodo release.

\begin{center}\rule{0.5\linewidth}{0.5pt}\end{center}

\emph{Linux Claude, 2026-04-15, under Yifan Chen's Gate 2 authorization.
Independent review of this appendix is archived as
\texttt{REVIEW\_PHASE\_B\_EXP1\_appendix.md}.}

\end{document}
